\chapter{Production Deployment Configuration}
\label{app:deployment}

This appendix summarizes the production deployment of \sh{} at \url{https://scholarhub.space}. The deployment is orchestrated by Docker Compose; the production compose file lives at the repository root. Sensitive values (passwords, API keys, signing secrets) are passed through environment variables and are not reproduced here.

\section{Service Topology}

The deployment consists of seven services, summarized in chapter \ref{ch:architecture} (\cref{tab:services}). The services run on a single host with the following resource allocation:

\begin{itemize}
    \item \textbf{nginx}: minimal CPU and memory footprint; one replica.
    \item \textbf{frontend}: static asset server; minimal footprint.
    \item \textbf{backend}: bounded by LLM and database I/O; one replica with async worker concurrency.
    \item \textbf{collab}: bounded by WebSocket connection count; one replica with Y.js room state in memory.
    \item \textbf{postgres}: persistent data store; volume-backed.
    \item \textbf{redis}: in-memory cache; volume-backed for AOF persistence.
    \item \textbf{certbot}: scheduled task for TLS certificate renewal.
\end{itemize}

\section{Environment Variables}

The principal environment variables required by the backend are listed in \cref{tab:envvars}. Values shown are illustrative; production values are specific to the deployment.

\small
\begin{longtable}{@{}>{\raggedright\arraybackslash\ttfamily\small}p{6.0cm} >{\raggedright\arraybackslash}p{\dimexpr\textwidth-6.6cm\relax}@{}}
\caption{Principal backend environment variables.}
\label{tab:envvars}\\
\toprule
\normalfont\textbf{Variable} & \textbf{Purpose} \\
\midrule
\endfirsthead
\multicolumn{2}{@{}l}{\textit{Table~\thetable{} (continued from previous page)}}\\
\toprule
\normalfont\textbf{Variable} & \textbf{Purpose} \\
\midrule
\endhead
\bottomrule
\multicolumn{2}{r@{}}{\textit{continued on next page}} \\
\endfoot
\bottomrule
\endlastfoot
DATABASE\_URL                    & PostgreSQL connection string. \\
REDIS\_URL                       & Redis connection string. \\
SECRET\_KEY                      & Application secret for signing session cookies. \\
OPENROUTER\_KEY\_ENCRYPTION\_KEY & Fernet key for encrypting user-supplied OpenRouter keys at rest. \\
COLLAB\_JWT\_SECRET              & Shared secret between backend and Hocuspocus for JWT verification. \\
COLLAB\_WS\_URL                  & WebSocket URL for the collaboration server. \\
OPENAI\_API\_KEY                 & OpenAI API key for first-party AI features. \\
OPENROUTER\_API\_KEY             & Fallback OpenRouter key for users without their own. \\
SEMANTIC\_SCHOLAR\_API\_KEY      & Authenticated key for Semantic Scholar's higher rate limit. \\
CORE\_API\_KEY                   & API key for the CORE service. \\
GOOGLE\_CLIENT\_ID, GOOGLE\_CLIENT\_SECRET & OAuth credentials. \\
DAILY\_API\_KEY                  & Daily.co API key for video meetings. \\
RESEND\_API\_KEY                 & Email service key for transactional email. \\
COOKIE\_DOMAIN                   & Domain for session cookie scope (\code{scholarhub.space}). \\
DEBUG                            & Must be \code{false} in production. \\
LATEX\_WARMUP\_ON\_STARTUP       & Pre-compile LaTeX templates on container start. \\
\end{longtable}

\section{Nginx Routing}

The Nginx configuration in \file{nginx/conf.d/scholarhub.conf} routes by path:

\begin{itemize}
    \item \code{/} -- static frontend bundle.
    \item \code{/api/} -- proxy to \code{backend:8000}.
    \item \code{/collab/} -- proxy to \code{collab:3001} with WebSocket upgrade.
    \item \code{/.well-known/acme-challenge/} -- certbot HTTP challenge.
\end{itemize}

TLS is terminated at Nginx using Let's Encrypt certificates renewed every 12 hours by the certbot container.

\section{Persistence and Volumes}

\begin{itemize}
    \item \code{postgres\_data} -- PostgreSQL data directory; durable.
    \item \code{redis\_data} -- Redis AOF and RDB files; durable.
    \item \code{./uploads} -- user-uploaded files (PDFs and images); durable.
    \item \code{./certbot/conf} -- TLS certificates and Let's Encrypt account state; durable.
\end{itemize}

\section{Health and Monitoring}

Each service defines a Docker health check. The backend exposes \code{/health} (basic) and \code{/healthz} (database and Redis verified). The collaboration server exposes a Node.js health check script. PostgreSQL uses \code{pg\_isready}; Redis uses \code{redis-cli ping}.

Logging is via Docker's default JSON file driver. Centralized log aggregation (Loki, Elasticsearch, or a hosted equivalent) is part of the productization plan in section \ref{sec:fw-multitenancy}.

\section{Disaster Recovery}

The current deployment relies on three recovery mechanisms:

\begin{enumerate}
    \item PostgreSQL data is backed up via filesystem snapshots of the Docker volume on a daily schedule.
    \item Redis data is persisted to disk through AOF; on container restart, state is reconstructed.
    \item User-uploaded PDFs in \code{./uploads} are part of the daily volume snapshot.
\end{enumerate}

A full disaster recovery test, including bare-metal restore from snapshots into a fresh deployment, has been performed during development to validate the procedure but is not part of a regular operational schedule. Productizing the system would require formalization of the recovery time objective (RTO) and recovery point objective (RPO) to satisfy institutional service level expectations.
